\section{Related Work}
\label{sec:related}

\subsection{Generative Inpainting and Diffusion Forensics}
Early image manipulation detection focused primarily on conventional digital forgeries, such as copy-move, splicing, and removal via patch-based texture synthesis~\cite{mvss,pscc}. The rise of deep generative models shifted research toward detecting GAN-generated imagery, where convolutional artifacts, spectral peaks in the frequency domain, and color-channel discrepancies served as reliable forensic indicators~\cite{cnndet,npr}. 

With the emergence of Latent Diffusion Models (LDMs)~\cite{ldm} and modern flow-matching architectures~\cite{flux}, generative inpainting has achieved unprecedented realism. Diffusion models generate images through iterative denoising score matching, introducing fine-grained stochastic noise variations rather than periodic frequency grid artifacts. To benchmark detection under these advanced generators, Wang et al.~\cite{opensdi} introduced the OpenSDID dataset, evaluating both image-level classification and pixel-level localization across multiple diffusion backbones. Recently, FUSED~\cite{fused} established benchmark results under an SD1.5-only training protocol using a massive Sparse Mixture-of-Experts (MoE) backbone on paired H100 GPUs. In this work, we address cross-generator generalization and local edit localization under a unified framework evaluated across diverse diffusion families.

\subsection{Foundation Vision Models in Forensics}
Foundation vision models pretrained on internet-scale data have demonstrated strong zero-shot transfer capabilities across vision tasks. In image forensics, Contrastive Language-Image Pretraining (CLIP) has been widely adopted; for instance, DeCLIP~\cite{declip} decodes intermediate CLIP representations to localize deepfake boundaries, while MaskCLIP~\cite{opensdi} combines CLIP semantic features with Masked Autoencoder (MAE) representations. 

However, CLIP-based forensic models inherit critical limitations. Because CLIP is trained via image-text contrastive objectives on aggressively downsampled square crops (typically $224 \times 224$), its feature representations prioritize high-level semantic alignment while actively suppressing sub-pixel high-frequency variations. Consequently, CLIP-based models struggle to resolve fine inpainting boundaries, particularly in small edits ($<2\%$ image area) where semantic content remains largely undisturbed. In contrast, self-supervised Vision Transformers such as DINOv2~\cite{dinov2} are trained with multi-crop masked image modeling without textual bottlenecks, producing spatially dense, feature-consistent patch tokens that delineate object boundaries with high geometric precision. In FRAME, we harness frozen DINOv2-Base patch tokens as a core semantic foundation.

\subsection{Forensic Noise Residual Extraction}
Digital image acquisition introduces characteristic sensor noise, color filter array (CFA) interpolation patterns, and compression artifacts that are systematically disrupted when synthetic pixels are inpainted into a camera photograph. Classical forensic methods extract these artifacts using Spatial Rich Models (SRM)~\cite{srm}, which apply fixed high-pass linear filters to suppress scene content and expose residual noise. 

To adapt residual extraction to learned representations, Bayar and Stamm~\cite{bayar} formulated constrained convolutional layers, which enforce zero-sum constraints on filter weights to guarantee that the central pixel is predicted from its local spatial neighborhood, thereby suppressing low-frequency semantic content. Modern forensic architectures such as TruFor~\cite{trufor} incorporate learned Noiseprint++ streams, while CAT-Net~\cite{catnet} mines DCT-domain compression traces. In this work, we construct a 60-channel forensic residual pyramid combining 30 fixed SRM filters and 30 learnable Bayar constrained convolutions, providing multi-scale noise anomaly evidence directly to our cross-attention decoder.

\subsection{Multi-Scale Localization and Supervision Paradigms}
Localizing inpainting boundaries requires balancing macroscopic scene coherence with microscopic boundary inspection. Existing methods such as MVSS-Net~\cite{mvss} and PSCC-Net~\cite{pscc} employ multi-scale progressive feature networks to segment spliced boundaries. However, when applied to generative inpainting, existing models exhibit severe performance degradation on tiny edits, where predicted IoU frequently collapses below $10\%$. 

Furthermore, historical training regimes have suffered from a critical supervision mismatch: models trained on local crops were supervised using global image-level labels~\cite{trufor,mvss}. When a crop extracted from an edited image fell entirely outside the manipulated region, the network was penalized for predicting an authentic label, inducing gradient conflicts and driving genuine-image false alarm rates to unacceptable levels. FRAME eliminates this supervision noise through a joint dual-view loss formulation that decouples whole-image classification from local crop presence, ensuring robust discrimination on both pristine and manipulated imagery.
